\slajdid{03-s014}
\begin{frame}[label=03-s014]
\fontsize{9}{10}\selectfont
\setlength{\parskip}{2pt}
U nekim situacijama multiplekseri su komponente koje omogućuju još jednostavniju realizaciju kombinacionih mreža od dekodera. Na primer ako treba da realizujemo funkciju \(F=\bar B\bar A+BA\). Pre nego što prikažemo relizaciju ono što treba uočiti da je multiplekser principski
\begingroup
\def\DEoznsize{8}\def\DEoblik{.6}\def\DEdekmod{1}\def\DEcs{0}\def\DEprimer{0}\def\DEjedinica{.75cm}
\begin{columns}[c,onlytextwidth]
\begin{column}{.44\textwidth}\centering
% DEdekmod=0: eksplicitni invertori; 1: blok dekodera. DEcs selekcija bloka.
\providecommand{\DEjedinica}{1cm}
\providecommand{\DEoblik}{.7}
\providecommand{\DEoznsize}{9}
\begin{tikzpicture}[circuit logic US,x=\DEjedinica,y=\DEjedinica,font=\fontsize{\DEoznsize}{\DEoznsize}\selectfont,
 inv/.style={not gate US,draw,minimum width=9mm,minimum height=8mm,scale=\DEoblik},
 kapija/.style={and gate US,draw,minimum width=10mm,minimum height=10mm,scale=\DEoblik,rotate=-90}]
\ifnum\DEdekmod=0
 \foreach \n/\y in {1/3.2,0/2.3}{
  \node[inv] (i\n) at (0,\y) {};
  \node[inv] (j\n) at (1.15,\y) {};
  \draw[DEvod] (i\n.input)--++(-.4,0) node[left] {S\n};
  \draw[DEvod] (i\n.output)--(j\n.input);
  \coordinate (tap\n) at ($(i\n.output)!.5!(j\n.input)$);
  \draw[DEvod] (j\n.output)--(5.4,\y);
  \draw[DEvod] (tap\n)--++(0,-.35)--(5.4,\y-.35);
 }
\else
 \draw[DEvod] (-.6,1.5) rectangle (1.2,3.85);
 \node[align=center] at (.3,3.4) {Dekoder\\2/4};
 \foreach \n/\y in {3/3.0,2/2.66,1/2.33,0/2.0}{
  \node[anchor=east,inner sep=1pt] at (1.2,\y) {Y\n};
  \draw[DEvod] (1.2,\y)--(5.4,\y);
 }
 \foreach \n/\x/\y in {1/-.18/1.23,0/.55/.96}{
  \node[anchor=south,inner sep=1pt] at (\x,1.5) {S\n};
  \draw[DEvod] (\x,1.5)--(\x,\y)--(-1,\y) node[left] {S\n};
 }
 \ifnum\DEcs=1
  \node[anchor=west,inner sep=1pt] at (-.6,2.03) {CS};
  \draw[DEvod] (-.6,2.03)--(-1,2.03) node[left] {CS};
 \fi
\fi
\ifnum\DEdekmod=1\def\DEvrh{4.2}\def\DEkutija{3.98}\else\def\DEvrh{3.98}\def\DEkutija{3.65}\fi
\foreach \n/\x/\sy/\ay/\dy in {0/2.05/2.85/1.95/2.0,1/3.05/2.85/2.3/2.33,2/4.05/3.2/1.95/2.66,3/5.05/3.2/2.3/3.0}{
 \ifnum\DEdekmod=0
  \node[kapija,logic gate inputs=nnn] (g\n) at (\x,.6) {};
  \draw[DEvod] (g\n.input 3)--(g\n.input 3 |- 0,\sy);
  \draw[DEvod] (g\n.input 2)--(g\n.input 2 |- 0,\ay);
 \else
  \node[kapija,logic gate inputs=nn] (g\n) at (\x,.6) {};
  \draw[DEvod] (g\n.input 2)--(g\n.input 2 |- 0,\dy);
 \fi
 \draw[DEvod] (g\n.input 1)--(g\n.input 1 |- 0,\DEvrh) node[above] {I\n};
}
\node[or gate US,draw,logic gate inputs=nnnn,minimum width=10mm,minimum height=12mm,scale=\DEoblik,rotate=-90] (z) at (3.55,-.92) {};
\foreach \n/\pin/\y in {0/4/-.22,1/3/0,2/2/0,3/1/-.22}
 \draw[DEvod] (g\n.output)--(g\n.output |- 0,\y)--(z.input \pin |- 0,\y)--(z.input \pin);
\draw[DEvod] (z.output)--++(0,-.25) node[below] {Z};
\coordinate (DEdno) at ($(z.output)+(0,-.08)$);
\draw[DEvod,dashed] (-.75,0 |- DEdno) rectangle (5.6,\DEkutija);
\node[align=center] at (.15,.15) {MUX\\4/1};
\end{tikzpicture}

\end{column}
\begin{column}{.21\textwidth}\centering
% DEcs: dodatni CS; DEprimer: spoljašnji nivoi 1,0,0,1 i B/A/F.
\begin{tikzpicture}[DEoznaka,x=1cm,y=1cm]
\ifnum\DEcs=1\def\DEpom{.35}\else\def\DEpom{0}\fi
\draw[DEvod] (0,0) rectangle (1.65,{2.7+\DEpom});
\node[align=center] at (.82,{2.25+\DEpom}) {MUX\\4/1};
\foreach \n/\y/\v in {3/1.72/1,2/1.32/0,1/.92/0,0/.52/1}{
 \draw[DEvod] (-.3,{\y+\DEpom})--(0,{\y+\DEpom});
 \node[anchor=west,inner sep=2pt] at (0,{\y+\DEpom}) {I\n};
 \ifnum\DEprimer=1
  \node[anchor=east] at (-.3,{\y+\DEpom}) {\v};
 \fi
}
\draw[DEvod] (1.65,{1.12+\DEpom})--(1.95,{1.12+\DEpom});
\node[anchor=east,inner sep=2pt] at (1.65,{1.12+\DEpom}) {Z};
\ifnum\DEprimer=1
 \node[anchor=west] at (1.95,{1.12+\DEpom}) {F};
\fi
\foreach \n/\x/\lab in {1/.48/B,0/1.18/A}{
 \draw[DEvod] (\x,0)--(\x,-.3);
 \node[anchor=south,inner sep=2pt] at (\x,0) {S\n};
 \ifnum\DEprimer=1
  \node[anchor=north] at (\x,-.3) {\lab};
 \fi
}
\ifnum\DEcs=1
 \draw[DEvod] (-.3,.48)--(0,.48);
 \node[anchor=west,inner sep=1pt] at (0,.48) {CS};
\fi
\end{tikzpicture}

\end{column}
\begin{column}{.06\textwidth}\centering
\begin{tikzpicture}
\draw[DEstrelica,DEplava] (0,0)--(.6,0);
\draw[DEstrelica,DEplava] (0,1.6)--(.6,.9);
\end{tikzpicture}
\end{column}
\begin{column}{.26\textwidth}\centering
\def\DEprimer{1}
% DEcs: dodatni CS; DEprimer: spoljašnji nivoi 1,0,0,1 i B/A/F.
\begin{tikzpicture}[DEoznaka,x=1cm,y=1cm]
\ifnum\DEcs=1\def\DEpom{.35}\else\def\DEpom{0}\fi
\draw[DEvod] (0,0) rectangle (1.65,{2.7+\DEpom});
\node[align=center] at (.82,{2.25+\DEpom}) {MUX\\4/1};
\foreach \n/\y/\v in {3/1.72/1,2/1.32/0,1/.92/0,0/.52/1}{
 \draw[DEvod] (-.3,{\y+\DEpom})--(0,{\y+\DEpom});
 \node[anchor=west,inner sep=2pt] at (0,{\y+\DEpom}) {I\n};
 \ifnum\DEprimer=1
  \node[anchor=east] at (-.3,{\y+\DEpom}) {\v};
 \fi
}
\draw[DEvod] (1.65,{1.12+\DEpom})--(1.95,{1.12+\DEpom});
\node[anchor=east,inner sep=2pt] at (1.65,{1.12+\DEpom}) {Z};
\ifnum\DEprimer=1
 \node[anchor=west] at (1.95,{1.12+\DEpom}) {F};
\fi
\foreach \n/\x/\lab in {1/.48/B,0/1.18/A}{
 \draw[DEvod] (\x,0)--(\x,-.3);
 \node[anchor=south,inner sep=2pt] at (\x,0) {S\n};
 \ifnum\DEprimer=1
  \node[anchor=north] at (\x,-.3) {\lab};
 \fi
}
\ifnum\DEcs=1
 \draw[DEvod] (-.3,.48)--(0,.48);
 \node[anchor=west,inner sep=1pt] at (0,.48) {CS};
\fi
\end{tikzpicture}

\end{column}
\end{columns}
\endgroup
odnosno da ponovo imamo selekciju potpunih proizvoda koje formiraju signali S1 i S0. U tom smislu jedino što treba da uradimo jeste da na ulaze multipleksera za odgovarajući proizvod dovedemo željene vrednosti funkcije
\end{frame}
